% ============================================================
% A system-level matrix formalism for coherent geometric dosimetry
% Companion note to monograph_v2.tex, Part III
% Author: R. Wydaeghe
% ============================================================

\documentclass[11pt,a4paper]{article}

\usepackage[margin=2.4cm]{geometry}
\usepackage[utf8]{inputenc}
\usepackage[T1]{fontenc}
\usepackage{lmodern}
\usepackage{microtype}
\linespread{1.04}

\usepackage{amsmath,amssymb,amsthm,mathtools}
\usepackage{bm}
\usepackage{mathrsfs}
\usepackage{bbm}
\usepackage{cancel}
\usepackage{xfrac}
\allowdisplaybreaks[2]

\usepackage{empheq}
\usepackage[most]{tcolorbox}
\tcbset{highlight math style={%
    enhanced, colframe=black!70, colback=gray!6,
    arc=2pt, boxrule=0.5pt, left=3pt, right=3pt}}

\usepackage{tikz}
\usetikzlibrary{arrows.meta,positioning,fit,backgrounds,calc}

\usepackage{enumitem}
\usepackage{xcolor}
\usepackage{xspace}
\usepackage{fancyhdr}
\pagestyle{fancy}
\fancyhf{}
\cfoot{\thepage}
\renewcommand{\headrulewidth}{0pt}

\usepackage{titlesec}
\titleformat{\section}{\large\bfseries}{\thesection}{1em}{}
\titleformat{\subsection}{\normalsize\bfseries}{\thesubsection}{1em}{}

\usepackage[colorlinks=true, linkcolor={blue!55!black},
            citecolor={green!45!black}, urlcolor={blue!65!black}]{hyperref}
\usepackage[nameinlink,capitalize]{cleveref}

% Cross-reference labels in monograph_v2.aux (same directory).
% xr-hyper preserves hyperref links across documents.
\usepackage{xr-hyper}
\IfFileExists{monograph_v2.aux}{\externaldocument{monograph_v2}}{}

\usepackage{aliascnt}
\theoremstyle{definition}
\newtheorem{theorem}{Theorem}[section]
\newaliascnt{proposition}{theorem}
\newtheorem{proposition}[proposition]{Proposition}
\aliascntresetthe{proposition}
\newaliascnt{lemma}{theorem}
\newtheorem{lemma}[lemma]{Lemma}
\aliascntresetthe{lemma}
\newaliascnt{corollary}{theorem}
\newtheorem{corollary}[corollary]{Corollary}
\aliascntresetthe{corollary}
\newaliascnt{definition}{theorem}
\newtheorem{definition}[definition]{Definition}
\aliascntresetthe{definition}
\newaliascnt{remark}{theorem}
\newtheorem{remark}[remark]{Remark}
\aliascntresetthe{remark}
\crefname{proposition}{Proposition}{Propositions}
\Crefname{proposition}{Proposition}{Propositions}
\crefname{lemma}{Lemma}{Lemmas}\Crefname{lemma}{Lemma}{Lemmas}
\crefname{corollary}{Corollary}{Corollaries}\Crefname{corollary}{Corollary}{Corollaries}
\crefname{definition}{Definition}{Definitions}\Crefname{definition}{Definition}{Definitions}
\crefname{remark}{Remark}{Remarks}\Crefname{remark}{Remark}{Remarks}

% Macros aligned with the monograph
\newcommand{\khat}{\hat{\bm{k}}}
\newcommand{\nhat}{\hat{\bm{n}}}
\newcommand{\rr}{\mathbf{r}}
\newcommand{\EE}{\mathbf{E}}
\newcommand{\Sab}{S_{\mathrm{ab}}}
\newcommand{\diff}{\mathrm{d}}
\newcommand{\Complex}{\mathbb{C}}
\newcommand{\Reals}{\mathbb{R}}
\DeclareMathOperator*{\argmax}{arg\,max}

% Local macros for the system formalism
\newcommand{\Jmat}{\mathbf{J}}                 % dispatch element -> path
\newcommand{\Phimat}{\mathbf{\Phi}}            % per-path phase at r
\newcommand{\Psimat}{\bm{\Psi}}                % stacked path polarisations
\newcommand{\Psitil}{\tilde{\bm{\Psi}}}        % tissue-weighted per-path readout
\newcommand{\Mmat}{\mathbf{M}}                 % path-path exposure Gram
\newcommand{\Qmat}{\mathbf{Q}}                 % element-element exposure operator
\newcommand{\Pmat}{\tilde{\mathbf{P}}}         % element-block projector J J^T
\newcommand{\xx}{\mathbf{x}}
\newcommand{\hh}{\mathbf{h}}
\newcommand{\cvec}{\mathbf{c}}
\newcommand{\bvec}{\mathbf{b}}
\newcommand{\alphavec}{\bm{\alpha}}
\newcommand{\psin}{\bm{\psi}_n}
\newcommand{\Fn}{\mathbf{F}_n}
\newcommand{\FnR}{\mathbf{F}_n(\rr)}
\newcommand{\GG}{\mathbf{G}}
\newcommand{\Gtil}{\tilde{\mathbf{G}}}
\newcommand{\Id}{\mathbf{I}}
\newcommand{\Sigmabody}{\Sigma}
\newcommand{\vv}{\mathbf{v}}
\newcommand{\uu}{\mathbf{u}}

\title{\textbf{A system-level matrix formalism\\ for coherent geometric dosimetry}\\[0.3em]
\large Companion note to \emph{Geometric Dosimetry}, Part~III}
\author{Robin Wydaeghe \\[-0.2em] \small Ghent University \& imec}
\date{\today}

\begin{document}
\maketitle

\begin{abstract}
Part~III of the monograph presents the coherent pipeline as a sequence of physical
operations: ray tracing, superposition, Fresnel transmission, depth
integration.  This note recasts the same pipeline as a short chain of matrix
factorisations acting on a shared state vector.  The communication channel
$\hh$, the field channel $\GG(\rr)$, and the exposure channel
$\Gtil(\rr)$ are revealed as three parallel read-outs of one common
\emph{path-amplitude carrier}
$\cvec(\rr,\xx)=\Phimat(\rr)\Jmat\xx$, where
$\Jmat\in\{0,1\}^{N\times M}$ is the element-to-path dispatch matrix and
$\Phimat(\rr)$ is a diagonal unitary of per-path phases at target point
$\rr$.  Once this split is made, the key results of Part~III admit one-line
proofs.  In particular, the exposure operator factors as
$\Qmat=\Jmat^T\Mmat\Jmat$ with $\Mmat\in\Complex^{N\times N}$ a
path-correlation Gram obtained from a single surface integral; the
worst-case absorbed power per unit transmit power equals
$\lambda_{\max}(\Mmat\Jmat\Jmat^T)$, an eigenvalue on the $N$-dimensional
path space rather than the $M$-dimensional element space; and the optimal
exposure-constrained precoder admits a Woodbury push-through form
$\xx^\star\propto\Jmat^T(\nu\Id_N+\lambda\Mmat\Jmat\Jmat^T)^{-1}\bvec$
that is diagonalised jointly with the path-space Gram.  The formalism
makes the multi-user, multi-body, and near-field extensions immediate, and
exposes the element-tied subspace $\operatorname{range}(\Jmat)\subset\Complex^N$
as the central geometric object separating what a precoder can and cannot
control.
\end{abstract}

% ============================================================
\section{Setup and motivation}\label{sec:motivation}
% ============================================================

The coherent pipeline of Part~III can be read as a plumbing diagram.  From
the base station, a precoding vector $\xx\in\Complex^M$ is broadcast into
the environment.  The ray tracer produces $N$ paths, each carrying a complex
polarisation-amplitude vector $\psin\in\Complex^3$ and a direction of
arrival $\khat_n$.  Three downstream observers read from the scene: the
physical electric field $\EE(\rr)$ in free space, the received signal $y$
at a user terminal at $\rr_\mathrm{UE}$, and the absorbed power density
$\Sab(\rr)$ on the body surface at $\rr\in\Sigmabody$.  Though the outputs
have different types (vector, scalar, scalar density), the machinery between
$\xx$ and any one of them is the same up to the last step.

What is unchanged across the three branches is the object
$x_{j(n)}\,e^{-ik_0\khat_n\cdot\rr}$: the beamforming weight of the element
that originated path $n$, phase-shifted by the path's geometric delay to
$\rr$.  What differs is only the per-path \emph{read-out} that follows:
$\psin$ itself for the field; $\mathbf{C}_R(\khat_n)^H\psin$ for the UE
signal; $\sqrt{\sigma/(4\alpha_n)}\FnR\psin$ for the tissue-coupled surface
field.  The monograph assembles these into three different matrices
$\GG,\hh,\Gtil$ that are then manipulated independently.  The present
note shows that there is a cheaper bookkeeping: factor out the common
carrier and let the branches live as three linear read-outs sitting in
parallel.

Once this is done, the quadratic dosimetry functional $\xx^H\Qmat\xx$,
the alignment $\rho$ between signal and exposure, and the optimal
exposure-constrained precoder all become statements about a single
$N\times N$ path-correlation Gram $\Mmat$ and its interaction with the
dispatch matrix $\Jmat$.  That compression is the motivation.

% ============================================================
\section{The path-amplitude carrier}\label{sec:carrier}
% ============================================================

Three objects organise the common plumbing.

\begin{definition}[Dispatch matrix]\label{def:J}
The \emph{element-to-path dispatch matrix} is
\begin{equation}\label{eq:J-def}
  \Jmat\in\{0,1\}^{N\times M},
  \qquad
  J_{nj}=\delta_{j(n),\,j}.
\end{equation}
Row $n$ has a single~$1$ in column $j(n)$, the transmit element that
originates path~$n$.  For $\xx\in\Complex^M$, the vector
$(\Jmat\xx)_n=x_{j(n)}$ broadcasts each element's weight onto all paths
leaving that element.
\end{definition}

\begin{definition}[Phase diagonal]\label{def:Phi}
For $\rr\in\Reals^3$ define the \emph{phase diagonal}
\begin{equation}\label{eq:Phi-def}
  \Phimat(\rr)=\operatorname{diag}\!\bigl(e^{-ik_0\khat_1\cdot\rr},\dots,
  e^{-ik_0\khat_N\cdot\rr}\bigr)\in\Complex^{N\times N}.
\end{equation}
$\Phimat(\rr)$ is a diagonal unitary; it applies to each path its
geometric phase at the observation point $\rr$.
\end{definition}

\begin{definition}[Path-amplitude carrier]\label{def:carrier}
The \emph{path-amplitude carrier} at observation point $\rr$ under
precoding~$\xx$ is the vector
\begin{equation}\label{eq:carrier}
  \boxed{\;
  \cvec(\rr,\xx)\;=\;\Phimat(\rr)\,\Jmat\,\xx\;\in\;\Complex^N.
  \;}
\end{equation}
Entry $n$ of $\cvec$ is the complex scalar
$c_n(\rr,\xx)=x_{j(n)}\,e^{-ik_0\khat_n\cdot\rr}$.
\end{definition}

$\cvec$ is the state variable from which every coherent observable follows.
Its space, $\Complex^N$, is the \emph{path space}; its complex structure
is entirely geometric (directions of arrival, path counts, phase centres),
while its magnitudes are entirely under the transmitter's control through
$\xx$.

Write $\Psimat=[\psi_1,\dots,\psi_N]\in\Complex^{3\times N}$ for the stack
of per-path polarisation-amplitude vectors.  The total electric field at
$\rr$ under $\xx$ is
\begin{equation}\label{eq:E-via-carrier}
  \EE(\rr)=\sum_n\psin\,c_n(\rr,\xx)=\Psimat\,\cvec(\rr,\xx)
         =\Psimat\,\Phimat(\rr)\,\Jmat\,\xx.
\end{equation}
This is exactly the field channel matrix
$\GG(\rr)=\Psimat\,\Phimat(\rr)\,\Jmat$ of \cref{eq:G-def} in the
monograph, now visibly factored into an environment-only block $\Psimat$,
a position-only block $\Phimat(\rr)$, and a combinatorial block $\Jmat$.

% ============================================================
\section{The three branches as parallel read-outs}\label{sec:branches}
% ============================================================

Three linear read-outs act on the shared carrier.

\paragraph{Field branch.} $\Psimat\in\Complex^{3\times N}$ acts on
$\cvec$ by left multiplication, returning the $\Complex^3$-valued
free-space field \eqref{eq:E-via-carrier}.  $\Psimat$ is determined by the
propagation environment alone (last-scatter geometry, transmit patterns);
it is position-independent.

\paragraph{Signal branch.} Collapse each path's polarisation to a scalar
by the UE antenna response, giving a row vector
\begin{equation}\label{eq:alpha-def}
  \alphavec^T=\bigl[\mathbf{C}_R(\khat_1)^H\psi_1,\dots,
  \mathbf{C}_R(\khat_N)^H\psi_N\bigr]\in\Complex^{1\times N}.
\end{equation}
The received signal at $\rr_\mathrm{UE}$ is
\begin{equation}\label{eq:y-via-carrier}
  y(\xx)=\alphavec^T\,\cvec(\rr_\mathrm{UE},\xx)
        =\alphavec^T\,\Phimat(\rr_\mathrm{UE})\,\Jmat\,\xx
        =\hh^T\xx,
\end{equation}
so the UE channel reads
$\hh^T=\alphavec^T\,\Phimat(\rr_\mathrm{UE})\,\Jmat$, or
equivalently
\begin{equation}\label{eq:h-path}
  \boxed{\;
  \hh^*\;=\;\Jmat^T\,\bvec,\qquad
  \bvec\equiv\Phimat(\rr_\mathrm{UE})^H\alphavec^*\in\Complex^N.
  \;}
\end{equation}
$\bvec$ is the path-space signal vector: its $n$-th entry is the complex
scalar contribution of path~$n$ to the UE after the antenna projection and
the UE phase.

\paragraph{Exposure branch.} For a surface point $\rr\in\Sigmabody$ with
normal $\nhat(\rr)$, define the \emph{tissue-weighted per-path read-out}
\begin{equation}\label{eq:psitil-def}
  \tilde{\bm\psi}_n(\rr)\;\equiv\;
  \sqrt{\tfrac{\sigma}{4\alpha_n}}\;\FnR\,\psin\in\Complex^3,
\end{equation}
with $\Fn$ the Fresnel transmission operator of \cref{eq:Fn} in the
monograph; stack them as
$\Psitil(\rr)=\bigl[\tilde{\bm\psi}_1(\rr),\dots,
\tilde{\bm\psi}_N(\rr)\bigr]\in\Complex^{3\times N}$.
The transmitted surface field is
\begin{equation}\label{eq:Etil-via-carrier}
  \tilde{\EE}(\rr)=\Psitil(\rr)\,\cvec(\rr,\xx)
                 =\Psitil(\rr)\,\Phimat(\rr)\,\Jmat\,\xx,
\end{equation}
and the absorbed power density is $\Sab(\rr)=\|\tilde{\EE}(\rr)\|^2$.
Thus $\Gtil(\rr)=\Psitil(\rr)\,\Phimat(\rr)\,\Jmat$, matching
\cref{eq:Gtilde-def}.

\paragraph{One picture, three read-outs.}  Taken together,
\Cref{fig:carrier} shows the system.  The operators differ only in the
last stage; all three share the dispatch block $\Jmat$ and the phase block
$\Phimat(\rr)$.

\begin{figure}[!h]
\centering
\begin{tikzpicture}[
  node distance=12mm and 16mm,
  box/.style={draw, rounded corners=2pt, minimum height=8mm, minimum width=16mm,
              align=center, inner sep=3pt, fill=gray!6},
  op/.style={draw, rounded corners=2pt, minimum height=7mm, minimum width=13mm,
             align=center, fill=white, font=\small},
  arr/.style={-{Stealth[length=2mm]}, thick},
]
  \node[box] (x) {$\xx\in\Complex^M$};
  \node[op, right=of x] (J) {$\Jmat$};
  \node[op, right=of J] (Phi) {$\Phimat(\rr)$};
  \node[box, right=of Phi] (c) {$\cvec(\rr,\xx)\in\Complex^N$};

  \node[op, above right=8mm and 10mm of c, fill=red!8] (alpha)
       {$\alphavec^T,\;\rr=\rr_\mathrm{UE}$};
  \node[op, right=of c, fill=gray!8] (PsiMat) {$\Psimat$};
  \node[op, below right=8mm and 10mm of c, fill=blue!8] (Psitilde)
       {$\Psitil(\rr),\;\rr\in\Sigmabody$};

  \node[right=10mm of alpha] (yout) {$y(\xx)\in\Complex$};
  \node[right=10mm of PsiMat] (Eout) {$\EE(\rr)\in\Complex^3$};
  \node[right=10mm of Psitilde] (Etout) {$\tilde\EE(\rr)\in\Complex^3$};
  \node[right=3mm of Etout, font=\small] (norm) {$\|\cdot\|^2\to\Sab(\rr)$};

  \draw[arr] (x) -- (J);
  \draw[arr] (J) -- (Phi);
  \draw[arr] (Phi) -- (c);
  \draw[arr] (c.east) -- ++(4mm,0) |- (alpha.west);
  \draw[arr] (c) -- (PsiMat);
  \draw[arr] (c.east) -- ++(4mm,0) |- (Psitilde.west);
  \draw[arr] (alpha) -- (yout);
  \draw[arr] (PsiMat) -- (Eout);
  \draw[arr] (Psitilde) -- (Etout);
  \draw[arr] (Etout) -- (norm);

  \begin{scope}[on background layer]
    \node[draw, dashed, rounded corners=2pt, fit=(J)(Phi)(c), inner sep=4pt,
          label={[font=\small\itshape, anchor=north]above:shared carrier}] {};
  \end{scope}
\end{tikzpicture}
\caption{The coherent pipeline as a two-stage factorisation.  The shared
carrier $\cvec(\rr,\xx)=\Phimat(\rr)\Jmat\xx$ is the state variable;
three read-outs feed off it in parallel.  The signal branch (top) collapses
to a scalar via the UE antenna row $\alphavec^T$ evaluated at
$\rr_\mathrm{UE}$; the field branch (middle) returns the free-space
$\Complex^3$ field via $\Psimat$; the exposure branch (bottom) returns the
tissue-coupled surface field via $\Psitil(\rr)$ and takes squared norm.}
\label{fig:carrier}
\end{figure}

% ============================================================
\section{Path-space factorisation of the exposure operator}
\label{sec:Q-factor}
% ============================================================

Integrating $\Sab$ over the body yields a quadratic form in $\xx$.  Using
the carrier factorisation, that form is pre-diagonalised.

\begin{definition}[Path-exposure Gram]\label{def:M}
The $N\times N$ Hermitian PSD \emph{path-exposure Gram} is
\begin{equation}\label{eq:M-def}
  \Mmat\;=\;\int_{\Sigmabody}
  \Phimat(\rr)^H\,\Psitil(\rr)^H\,\Psitil(\rr)\,\Phimat(\rr)\,\diff A,
  \qquad\Mmat\in\Complex^{N\times N}.
\end{equation}
Entry $(n,n')$ is
\begin{equation}\label{eq:M-entry}
  M_{nn'}\;=\;\sqrt{\tfrac{\sigma}{4\alpha_n}}
  \sqrt{\tfrac{\sigma}{4\alpha_{n'}}}
  \int_{\Sigmabody}
  \psi_n^H\,\Fn(\rr)^H\,\mathbf{F}_{n'}(\rr)\,\psi_{n'}\,
  e^{-ik_0(\khat_{n'}-\khat_n)\cdot\rr}\,\diff A.
\end{equation}
\end{definition}

\begin{theorem}[Path-space factorisation of the exposure operator]%
\label{thm:Q-factor}
The exposure operator of \cref{def:Q} in the monograph satisfies
\begin{equation}\label{eq:Q-JMJ}
  \boxed{\;
  \Qmat\;=\;\Jmat^T\,\Mmat\,\Jmat,
  \qquad P_\mathrm{abs}(\xx)=\xx^H\Qmat\xx.
  \;}
\end{equation}
\end{theorem}

\begin{proof}
From $\Gtil(\rr)=\Psitil(\rr)\Phimat(\rr)\Jmat$,
\[
  \Qmat=\int_{\Sigmabody}\Gtil(\rr)^H\Gtil(\rr)\,\diff A
       =\Jmat^T\!\int_{\Sigmabody}\Phimat(\rr)^H\Psitil(\rr)^H
        \Psitil(\rr)\Phimat(\rr)\,\diff A\;\Jmat
       =\Jmat^T\Mmat\Jmat.
\]
\end{proof}

The physical content is that $\Mmat$ is a path-indexed object: it knows
about pairs of paths, not about transmit elements.  The element structure
enters only through the aggregator $\Jmat$ in the last factorisation step.

\begin{remark}[Fourier content of $\Mmat$]\label{rem:fourier}
The phase kernel in \eqref{eq:M-entry} is
$e^{-i\mathbf q_{nn'}\cdot\rr}$ with
$\mathbf q_{nn'}\equiv k_0(\khat_{n'}-\khat_n)$.  The off-diagonal
entries of $\Mmat$ are therefore body-surface Fourier components at the
\emph{difference} wave-vector of each pair of paths, weighted by the
vector inner product of their Fresnel-filtered polarisations.  Diagonal
entries ($n=n'$) reduce to the incoherent self-absorption integrated over
$\Sigmabody$.  The full spectrum of $\Mmat$ is therefore controlled by
(a) the geometric pair-wise angular separation of paths, which sets
$|\mathbf q_{nn'}|$, and (b) how the body surface projects onto each
$\mathbf q_{nn'}$, which sets the magnitude of the Fourier component.
Flat patches with dimensions large compared to $2\pi/|\mathbf q_{nn'}|$
contribute little except near $\mathbf q_{nn'}\approx 0$.
\end{remark}

\begin{remark}[Element-block projector]\label{rem:P}
Let
\begin{equation}\label{eq:Ptilde}
  \Pmat\;\equiv\;\Jmat\Jmat^T\in\{0,1\}^{N\times N},
  \qquad(\Pmat)_{nn'}=\delta_{j(n),\,j(n')}.
\end{equation}
$\Pmat$ is block-diagonal in a path-ordering that groups co-originating
paths: each block is an all-ones matrix of size $N_j\times N_j$.  With the
normalised dispatch $\hat\Jmat\equiv\Jmat\operatorname{diag}(1/\sqrt{N_j})$,
$\hat\Jmat^T\hat\Jmat=\Id_M$ and
$\hat\Jmat\hat\Jmat^T=\Jmat\operatorname{diag}(1/N_j)\Jmat^T$
is an orthogonal projector onto
\begin{equation}\label{eq:element-tied}
  \mathcal E\;\equiv\;\operatorname{range}(\Jmat)
  =\{\mathbf v\in\Complex^N:v_n=v_{n'}\text{ whenever }j(n)=j(n')\}
  \subset\Complex^N.
\end{equation}
$\mathcal E$ is the \emph{element-tied subspace}: the set of path-space
vectors whose co-originating entries are equal.  It is the space of
degrees of freedom a precoder can actually excite.
\end{remark}

% ============================================================
\section{Worst-case absorption as a path-space eigenvalue}
\label{sec:wc-eig}
% ============================================================

\Cref{thm:Q-factor} reduces the spectral analysis of $\Qmat$ to that of
an $N\times N$ path-space object.

\begin{corollary}[Worst-case eigenvalue in path space]\label{cor:lam-path}
The nonzero eigenvalues of $\Qmat=\Jmat^T\Mmat\Jmat$ coincide with those
of $\Mmat\Pmat=\Mmat\Jmat\Jmat^T$; in particular,
\begin{equation}\label{eq:lam-path}
  \lambda_{\max}(\Qmat)\;=\;\lambda_{\max}(\Mmat\,\Pmat).
\end{equation}
A worst-case precoder direction $\xx_\star$ lifts from the path-space
eigenvector $\mathbf v_\star\in\mathcal E$ of $\Mmat\Pmat$ via
$\xx_\star\propto\Jmat^T\mathbf v_\star$.
\end{corollary}

\begin{proof}
For any matrices $A\in\Complex^{p\times q}$, $B\in\Complex^{q\times p}$,
the nonzero spectra of $AB$ and $BA$ coincide.  Apply with
$A=\Jmat^T$, $B=\Mmat\Jmat$: nonzero eigenvalues of
$\Jmat^T(\Mmat\Jmat)=\Qmat$ equal those of
$(\Mmat\Jmat)\Jmat^T=\Mmat\Pmat$.  The eigenvector lift follows from
the standard push-through:
$\Qmat\xx_\star=\lambda_\star\xx_\star$ with
$\xx_\star=\Jmat^T\mathbf v_\star$ whenever
$\Mmat\Pmat\mathbf v_\star=\lambda_\star\mathbf v_\star$ and
$\mathbf v_\star\in\mathcal E$.
\end{proof}

The picture is that $\Qmat$ lives on $\Complex^M$, but only the
$\le\operatorname{rank}(\Jmat)=M$ dimensions of the element-tied subspace
$\mathcal E\subset\Complex^N$ carry its nonzero spectrum.  Absorption
modes of the body are path-space vectors restricted to~$\mathcal E$; the
remainder of path space (the $N-M$ directions orthogonal to~$\mathcal E$)
is unreachable by any precoder.

% ============================================================
\section{Signal--exposure alignment in path space}\label{sec:rho-path}
% ============================================================

The alignment parameter $\rho$ of \cref{eq:rho-def} in the monograph
expresses, in element-space, how much of the MRT direction $\hh^*$ points
along top eigenmodes of $\Qmat$.  In path space it becomes transparent.

\begin{proposition}[Path-space alignment]\label{prop:rho-path}
With $\hh^*=\Jmat^T\bvec$ as in \eqref{eq:h-path},
\begin{equation}\label{eq:rho-path}
  \rho\;=\;\frac{\bvec^H\,\Pmat\,\Mmat\,\Pmat\,\bvec}
  {\bigl(\bvec^H\Pmat\bvec\bigr)\,\lambda_{\max}(\Mmat\Pmat)},
  \qquad \rho\in[0,1].
\end{equation}
\end{proposition}

\begin{proof}
Since $\Jmat$ has real entries and $\Phimat_U\equiv\Phimat(\rr_\mathrm{UE})$
is diagonal, $\hh^*=\Jmat^T\bvec$ with $\bvec=\Phimat_U^H\alphavec^*$ and
hence $\hh=\Jmat^T\bvec^*$, so $\hh^T=\bvec^H\Jmat$.  Then
\begin{align*}
  \|\hh\|^2 &\;=\;\|\hh^*\|^2
           \;=\;(\Jmat^T\bvec)^H(\Jmat^T\bvec)
           \;=\;\bvec^H\Jmat\Jmat^T\bvec
           \;=\;\bvec^H\Pmat\bvec,\\[0.3em]
  \hh^T\Qmat\hh^* &\;=\;(\bvec^H\Jmat)\bigl(\Jmat^T\Mmat\Jmat\bigr)(\Jmat^T\bvec)
                  \;=\;\bvec^H(\Jmat\Jmat^T)\Mmat(\Jmat\Jmat^T)\bvec
                  \;=\;\bvec^H\Pmat\Mmat\Pmat\bvec.
\end{align*}
Substituting into \cref{eq:rho-def} of the monograph together with
$\lambda_{\max}(\Qmat)=\lambda_{\max}(\Mmat\Pmat)$ (\Cref{cor:lam-path})
gives \eqref{eq:rho-path}.  The upper bound $\rho\le 1$ follows from the
element-space Rayleigh bound in \cref{eq:rho-def}; equivalently, setting
$\xx=\Jmat^T\bvec=\hh^*$ and using
$\xx^H\Qmat\xx\le\lambda_{\max}(\Qmat)\|\xx\|^2$.
\end{proof}

\begin{remark}[Reading $\rho$ geometrically]\label{rem:rho-geom}
Three path-space ingredients control $\rho$:
\begin{enumerate}[leftmargin=1.7em,itemsep=0.2em]
  \item the path-space signal vector $\bvec$ (which paths contribute to
  signal and with what phase);
  \item the element-block projector $\Pmat$ (how paths are grouped by
  transmit element);
  \item the path-exposure Gram $\Mmat$ (how paths correlate on the body
  through the Fresnel filter and surface-Fourier kernel).
\end{enumerate}
$\rho\to 0$ when $\Pmat\bvec$ lies in the null space of $\Mmat$: the
signal-aligned paths, after element-grouping, do not illuminate the body
(grazing, back-facing, or surface-Fourier-destructive pairs).  $\rho=1$
when $\Pmat\bvec$ is the top eigenvector of $\Mmat\Pmat$.
\end{remark}

% ============================================================
\section{Exposure-constrained beamforming via Woodbury push-through}
\label{sec:ecbf-path}
% ============================================================

The closed-form optimal precoder of \cref{eq:optimal-x} in the monograph,
\[
  \xx^\star=\sqrt P\,\frac{(\lambda\Qmat+\nu\Id_M)^{-1}\hh^*}
  {\|(\lambda\Qmat+\nu\Id_M)^{-1}\hh^*\|},
\]
appears to require inverting an $M\times M$ matrix that mixes the full
path geometry through $\Qmat$.  In fact it lifts to a path-space inverse.

\begin{theorem}[Path-space optimal precoder]\label{thm:x-path}
Let $\bvec$ be as in \eqref{eq:h-path}.  Then
\begin{equation}\label{eq:x-path}
  \boxed{\;
  \xx^\star\;\propto\;\Jmat^T\,\bigl(\nu\,\Id_N+\lambda\,\Mmat\,\Pmat\bigr)^{-1}\bvec,
  \;}
\end{equation}
where the proportionality is the normalisation of \cref{eq:optimal-x}.
The relevant spectral problem therefore lives on $\Complex^N$
with operator $\Mmat\Pmat$, which has rank $\le M$.
\end{theorem}

\begin{proof}
Apply the push-through identity $(\Id_M+AB)^{-1}A=A(\Id_N+BA)^{-1}$ with
$A=\Jmat^T\in\Complex^{M\times N}$ and
$B=(\lambda/\nu)\Mmat\Jmat\in\Complex^{N\times M}$.  Then
$AB=(\lambda/\nu)\Jmat^T\Mmat\Jmat=(\lambda/\nu)\Qmat$ and
$BA=(\lambda/\nu)\Mmat\Pmat$, so
\[
  (\Id_M+\tfrac{\lambda}{\nu}\Qmat)^{-1}\Jmat^T
  =\Jmat^T(\Id_N+\tfrac{\lambda}{\nu}\Mmat\Pmat)^{-1}.
\]
Multiplying by $1/\nu$ and applying to $\bvec$ gives
$(\lambda\Qmat+\nu\Id_M)^{-1}\hh^*
=\Jmat^T(\nu\Id_N+\lambda\Mmat\Pmat)^{-1}\bvec$.
\end{proof}

\begin{remark}[Diagonalisation in the exposure eigenbasis]%
\label{rem:diag-M}
Let $\Qmat=\sum_k\mu_k\vv_k\vv_k^H$ be the (Hermitian) element-space
eigendecomposition of the exposure operator, with $\vv_k\in\Complex^M$
orthonormal and $\mu_1\ge\cdots\ge\mu_M\ge 0$.  By \Cref{cor:lam-path},
each nonzero $\mu_k$ lifts from a right eigenvector $\uu_k\in\Complex^N$
of the non-Hermitian operator $\Mmat\Pmat$ via
$\vv_k=\Jmat^T\uu_k/\sqrt{\uu_k^H\Pmat\uu_k}$.  The ECBF solution
\cref{eq:optimal-x} of the monograph diagonalises in the $\{\vv_k\}$ basis
as
\begin{equation}\label{eq:x-diag}
  \xx^\star\;\propto\;\sum_{k=1}^{M}\frac{\vv_k\vv_k^H\hh^*}{\nu+\lambda\mu_k}
  \;=\;\Jmat^T\sum_{k=1}^{M}\frac{\uu_k\,\uu_k^H\,\Pmat\,\bvec}
  {\bigl(\uu_k^H\Pmat\uu_k\bigr)\,(\nu+\lambda\mu_k)},
\end{equation}
where the second form uses $\hh^*=\Jmat^T\bvec$ and
$\Jmat\Jmat^T=\Pmat$.  The shaping is explicit: each mode is
attenuated by $1/(\nu+\lambda\mu_k)$, and modes aligned with
$\ker(\Mmat\Pmat)$ (the absorption null-space) are unattenuated.  In the
limit $\lambda\to 0$, \eqref{eq:x-diag} collapses to MRT
$\xx^\star\propto\hh^*=\Jmat^T\bvec$; as $\lambda\to\infty$, only modes
with $\mu_k=0$ survive, i.e.\ the precoder steers strictly through the
exposure null-space.
\end{remark}

\begin{corollary}[Complexity]\label{cor:complexity}
Solving \eqref{eq:QCQP} requires only the nonzero eigenpairs of $\Mmat\Pmat$.
Since $\operatorname{rank}(\Mmat\Pmat)\le M$, at most $M$ eigenpairs must
be computed.  Once done, the optimal precoder at any Lagrange pair
$(\lambda,\nu)$ follows from \eqref{eq:x-diag} in $O(MN)$ flops for the
$\Jmat^T$ lift and $O(M)$ scalar operations for the mode shaping.  No
$M\times M$ inverse is required at runtime.
\end{corollary}

The cost driver is therefore the off-line assembly and partial
eigendecomposition of $\Mmat\Pmat$, exactly where the body geometry and
path structure interact.

% ============================================================
\section{Multi-user and multi-body: stacked read-outs}\label{sec:MU}
% ============================================================

The carrier-and-read-outs picture extends without new physics.  With $K$
UEs indexed by $k$ and $U$ bodies indexed by $u$:

\begin{itemize}[leftmargin=1.7em,itemsep=0.2em]
  \item The dispatch $\Jmat$ and path family $\{\khat_n,\psin\}$ are
  environment-shared; $\Phimat(\rr)$ is position-shared.
  \item Each UE contributes its own signal read-out
  $(\alphavec^{(k)},\rr_\mathrm{UE}^{(k)})$, hence its own path-space
  vector $\bvec^{(k)}=\Phimat(\rr_\mathrm{UE}^{(k)})^H\alphavec^{(k)*}$ and
  channel $\hh_k^*=\Jmat^T\bvec^{(k)}$.
  \item Each body contributes its own exposure read-out
  $\Psitil^{(u)}(\rr)$ on its surface $\Sigmabody^{(u)}$, hence a
  per-body path-exposure Gram
  \begin{equation}\label{eq:Mu-def}
    \Mmat^{(u)}=\int_{\Sigmabody^{(u)}}
    \Phimat(\rr)^H\,\Psitil^{(u)}(\rr)^H\,
    \Psitil^{(u)}(\rr)\,\Phimat(\rr)\,\diff A,
  \end{equation}
  and a per-body exposure operator
  $\Qmat^{(u)}=\Jmat^T\Mmat^{(u)}\Jmat$.
\end{itemize}

The MU-MIMO exposure-constrained problem \cref{eq:MU-QCQP,eq:Pabs-MU} in
the monograph becomes, in path space,
\begin{equation}\label{eq:mu-path}
  \max_{\{\mathbf w_k\}}\;\sum_k
  \log_2\!\left(1+\frac{|\bvec^{(k)\,T}\Jmat\mathbf w_k|^2}
  {\sum_{k'\neq k}|\bvec^{(k)\,T}\Jmat\mathbf w_{k'}|^2+\sigma_n^2}\right)
\end{equation}
subject to $\sum_k\mathbf w_k^H\Jmat^T\Mmat^{(u)}\Jmat\mathbf w_k\le
P_\mathrm{abs}^{\max}$ for each $u$ and
$\sum_k\|\mathbf w_k\|^2\le P$.  The path-space lifts in $\mathbf w_k$
of \Cref{thm:x-path} extend to each $(\mathbf w_k)$ separately at each
WMMSE iterate; the per-body matrices $\Mmat^{(u)}$ enter as a family of
PSD constraints on the same path space.

\begin{remark}[Structural separation]\label{rem:MU-sep}
The variables that change with the scene (bodies move, UEs move, antennas
repoint) populate $\Mmat^{(u)}$ and $\bvec^{(k)}$.  The combinatorial
object $\Jmat$, which is usually the bulk of $N$, is determined by the
transmitter hardware and the ray-tracing environment; it is expected to
be stable over many channel realisations.  The formalism therefore admits
a clean caching split: recompute only $\Mmat^{(u)}$ and $\bvec^{(k)}$
between coherence intervals.
\end{remark}

% ============================================================
\section{Isotropic transmit and trace identities}\label{sec:trace}
% ============================================================

The factorisation $\Qmat=\Jmat^T\Mmat\Jmat$ links the diagonal spectral
quantity of $\Qmat$, its trace, to a compact path-space sum.

\begin{proposition}[Isotropic-transmit absorbed power]\label{prop:iso}
Let $\xx\in\Complex^M$ be uniformly distributed on the complex sphere
$\{\xx:\|\xx\|^2=P\}$.  Then
\begin{equation}\label{eq:iso}
  \mathbb E_\xx\bigl[P_\mathrm{abs}(\xx)\bigr]
  \;=\;\frac{P}{M}\,\operatorname{tr}(\Qmat)
  \;=\;\frac{P}{M}\,\operatorname{tr}(\Mmat\,\Pmat).
\end{equation}
\end{proposition}

\begin{proof}
Uniformity on the sphere gives $\mathbb E[\xx\xx^H]=(P/M)\Id_M$; hence
$\mathbb E[\xx^H\Qmat\xx]=\operatorname{tr}(\Qmat\,\mathbb E[\xx\xx^H])
=(P/M)\operatorname{tr}(\Qmat)$.  The second equality uses the cyclic
property and $\Qmat=\Jmat^T\Mmat\Jmat$:
$\operatorname{tr}(\Jmat^T\Mmat\Jmat)
=\operatorname{tr}(\Mmat\Jmat\Jmat^T)=\operatorname{tr}(\Mmat\Pmat)$.
\end{proof}

\begin{corollary}[Breakdown of the trace]\label{cor:tr-structure}
With $(\Pmat)_{nn'}=\delta_{j(n),j(n')}$ from \eqref{eq:Ptilde},
\begin{equation}\label{eq:tr-breakdown}
  \operatorname{tr}(\Mmat\Pmat)
  \;=\;\sum_{n}M_{nn}
  \;+\;\sum_{n\ne n'}\,\delta_{j(n),j(n')}\,M_{n'n},
\end{equation}
so the isotropic-transmit absorbed power decomposes into an incoherent
self-absorption sum (all paths) plus a coherent within-element
cross-correlation (paths leaving the same transmit element).  Paths from
distinct elements never contribute: their cross-terms are averaged out by
the uniform phase of independent precoder entries.
\end{corollary}

\begin{proof}
$\operatorname{tr}(\Mmat\Pmat)=\sum_{n,n'}M_{nn'}P_{n'n}
=\sum_{n,n'}M_{nn'}\,\delta_{j(n),j(n')}$.  Split into $n=n'$ and $n\ne n'$
terms.
\end{proof}

\Cref{cor:tr-structure} recovers a monograph result in one line: the
``incoherent limit'' of \cref{cor:incoherent} in Part~III equals the
diagonal of \eqref{eq:tr-breakdown}; the extra within-element coherent
cross-terms are the angular-diversity gain of a multi-path element.  The
same trace identity bounds the worst-case eigenvalue by
\begin{equation}\label{eq:lam-tr-bound}
  \lambda_{\max}(\Qmat)\le\operatorname{tr}(\Qmat)=\operatorname{tr}(\Mmat\Pmat),
\end{equation}
which is tight when $\Qmat$ is rank-one (one dominant absorption mode).

% ============================================================
\section{Broadband extension}\label{sec:broadband}
% ============================================================

All path quantities ($\alpha_n,\beta_n,t_{s,n},t_{p,n},\psin,\khat_n$) depend
on frequency through $k_0=\omega/c$ and the tissue refractive index
$\tilde{n}(\omega)$.  The factorisation of
\Cref{sec:carrier,sec:Q-factor} holds at each frequency independently.

\begin{proposition}[Frequency-pointwise factorisation]\label{prop:omega-pt}
Let $\xx(\omega)\in\Complex^M$ be a square-integrable precoder with
transmit-energy constraint
$\|\xx\|_{L^2}^2\equiv\int|\xx(\omega)|^2\,\diff\omega/(2\pi)\le P$.  Under
the standard assumption that distinct frequency bins carry uncorrelated
data, the time-averaged absorbed power is
\begin{equation}\label{eq:Pabs-omega}
  \bar{P}_\mathrm{abs}
  \;=\;\int \xx(\omega)^H\,\Qmat(\omega)\,\xx(\omega)\,\frac{\diff\omega}{2\pi},
  \qquad
  \Qmat(\omega)=\Jmat^T\Mmat(\omega)\Jmat,
\end{equation}
and the broadband ECBF problem
\begin{equation}\label{eq:ecbf-omega}
  \max_{\xx(\cdot)}\int\bigl|\hh(\omega)^T\xx(\omega)\bigr|^2\,\frac{\diff\omega}{2\pi}
  \quad\text{s.t.}\quad
  \bar{P}_\mathrm{abs}\le P_\mathrm{abs}^{\max},\;\|\xx\|_{L^2}^2\le P
\end{equation}
decouples pointwise: the optimal precoder satisfies, at each $\omega$,
\begin{equation}\label{eq:x-omega}
  \xx^\star(\omega)\;\propto\;\Jmat^T\bigl(\nu\,\Id_N
  +\lambda\,\Mmat(\omega)\Pmat\bigr)^{-1}\bvec(\omega),
\end{equation}
with scalar multipliers $\lambda,\nu\ge 0$ shared across all frequencies and
determined by the two integrated constraints.
\end{proposition}

\begin{proof}
Stream independence gives
$\mathbb E\bigl[\int|\EE^\mathrm{trans}(\rr,t)|^2\diff t\bigr]
=\int\Sab(\rr,\omega)\diff\omega/(2\pi)$ by Parseval, because cross-$(\omega,\omega')$
terms in $|\sum_\omega\EE^\mathrm{trans}|^2$ integrate to zero under uncorrelated
symbols.  \Cref{thm:coherent-law} of the monograph applies to each
$\omega$-slice, giving \eqref{eq:Pabs-omega}.  The Lagrangian of
\eqref{eq:ecbf-omega} separates over $\omega$:
$\mathcal L=\int[\,|\hh^T\xx|^2-\lambda\xx^H\Qmat\xx-\nu|\xx|^2\,]\diff\omega/(2\pi)
+\lambda P_\mathrm{abs}^{\max}+\nu P$, and the stationary condition at each
$\omega$ is \cref{eq:optimal-x} of the monograph with the path-space
push-through of \Cref{thm:x-path}.
\end{proof}

\begin{remark}[What is not decoupled]\label{rem:omega-coupling}
The integrated constraints couple the frequencies through the shared
duals $\lambda,\nu$: more absorption budget at one frequency means less at
another.  Per-frequency eigenspectra
$\{\mu_k(\omega)\}_{k=1}^M$ of $\Mmat(\omega)\Pmat$ therefore jointly
determine the power-allocation law, analogous to water-filling in
capacity problems but with the ``cost'' set by $\mu_k(\omega)$ rather than
inverse noise.  The frequency-flat special case recovers
\Cref{thm:x-path}.
\end{remark}

% ============================================================
\section{Near-field and non-unitary phase kernels}\label{sec:nearfield}
% ============================================================

The locally-plane-wave assumption $d_n>3\lambda$ (distance from path's
last scatter to body) holds for most outdoor and many indoor
configurations but fails for on-body or close-proximity sources.  The
formalism carries over with one change.

\begin{definition}[Near-field phase kernel]\label{def:NF}
Let $\rr_n^{(\mathrm{scat})}$ be the last-scatter point of path~$n$,
$R_n(\rr)=|\rr-\rr_n^{(\mathrm{scat})}|$ the distance from it to the
observation point, and fix any reference point $\rr^{(0)}$ on the body.
Define the \emph{near-field phase kernel}
\begin{equation}\label{eq:PhiNF}
  \Phimat^{\mathrm{NF}}(\rr)=\operatorname{diag}\bigl(\varphi_n(\rr)\bigr),
  \qquad
  \varphi_n(\rr)=\frac{R_n^{(0)}}{R_n(\rr)}\,
  e^{-ik_0\bigl(R_n(\rr)-R_n^{(0)}\bigr)},
\end{equation}
with $R_n^{(0)}\equiv R_n(\rr^{(0)})$.  By construction $\varphi_n(\rr^{(0)})=1$
and $\psin$ is redefined to absorb the plane-wave-equivalent amplitude and
absolute phase at $\rr^{(0)}$ (as in \cref{sec:antenna-paths} of the
monograph, evaluated at $\rr^{(0)}$ rather than at infinity).
\end{definition}

Unlike the plane-wave $\Phimat(\rr)$, $\Phimat^{\mathrm{NF}}(\rr)$ is
\emph{diagonal but not unitary}: its entries carry both a phase
$e^{-ik_0 R_n}$ and a $1/R_n$ amplitude decay.  The Fresnel operator
$\Fn$ now also depends on $\rr$ through the direction
$\khat_n(\rr)=(\rr-\rr_n^{(\mathrm{scat})})/R_n(\rr)$, which varies over
the body.

\begin{theorem}[Near-field carrier factorisation]\label{thm:NF}
With $\Phimat$ replaced by $\Phimat^{\mathrm{NF}}$ and $\Fn$ by
$\Fn(\khat_n(\rr))$,
\begin{equation}\label{eq:NF-factor}
  \GG^{\mathrm{NF}}(\rr)=\Psimat\,\Phimat^{\mathrm{NF}}(\rr)\,\Jmat,
  \qquad
  \Gtil^{\mathrm{NF}}(\rr)=\Psitil^{\mathrm{NF}}(\rr)\,\Phimat^{\mathrm{NF}}(\rr)\,\Jmat,
\end{equation}
so the shared carrier still reads
$\cvec^{\mathrm{NF}}(\rr,\xx)=\Phimat^{\mathrm{NF}}(\rr)\Jmat\xx$.  The
path-exposure Gram takes the same form
\begin{equation}\label{eq:M-NF}
  \Mmat^{\mathrm{NF}}=\int_{\Sigmabody}
  \Phimat^{\mathrm{NF}}(\rr)^H\,\Psitil^{\mathrm{NF}}(\rr)^H\,
  \Psitil^{\mathrm{NF}}(\rr)\,\Phimat^{\mathrm{NF}}(\rr)\,\diff A,
\end{equation}
and $\Qmat=\Jmat^T\Mmat^{\mathrm{NF}}\Jmat$,
$\lambda_{\max}(\Qmat)=\lambda_{\max}(\Mmat^{\mathrm{NF}}\Pmat)$,
\Cref{thm:x-path} all proceed verbatim.
\end{theorem}

\begin{proof}
The only place the plane-wave assumption enters the derivation is in the
phase factor $e^{-ik_0\khat_n\cdot\rr}$ and in the $\rr$-independence of
$\khat_n$.  Replacing both by their near-field counterparts preserves the
structure $\GG=\Psimat\Phimat\Jmat$, $\Gtil=\Psitil\Phimat\Jmat$; the
push-through \Cref{thm:x-path} used only $\Qmat=\Jmat^T\Mmat\Jmat$ and
did not rely on the specific kernel of $\Mmat$.
\end{proof}

\begin{remark}[Confocal Fourier kernel]\label{rem:confocal}
Off-diagonal entries of $\Mmat^{\mathrm{NF}}$ contain the kernel
$\overline{\varphi_n(\rr)}\varphi_{n'}(\rr)
=\tfrac{R_n^{(0)}R_{n'}^{(0)}}{R_n(\rr)R_{n'}(\rr)}\,
e^{-ik_0(R_{n'}(\rr)-R_n(\rr))}$ times a constant pair-phase, whose
oscillatory factor has phase level-sets of constant path-length
difference: surfaces that are confocal ellipsoids with foci at the two
last-scatter points $\rr_n^{(\mathrm{scat})},\rr_{n'}^{(\mathrm{scat})}$.
Whereas the far-field kernel $e^{-i\mathbf q_{nn'}\cdot\rr}$ oscillates
at a constant wave-vector, the near-field kernel oscillates at a
position-dependent effective wave-vector and carries an $R^{-1}$-type
amplitude envelope.  Coherent hotspot shapes shift from tight Fourier
resolution cells toward confocal lobes.
\end{remark}

% ============================================================
\section{Stochastic channels and expected exposure}\label{sec:stochastic}
% ============================================================

Ray tracing gives a deterministic snapshot, but real scenes fluctuate:
moving scatterers, small body-pose changes, antenna calibration drift,
and phase-noise all randomise per-path parameters.  The formalism
separates cleanly.

Let $\Mmat$ be a random $N\times N$ Hermitian PSD matrix induced by
stochastic $\{\psin,\khat_n\}$, with finite-mean, finite-second-moment
distribution.  Write $\bar\Mmat\equiv\mathbb E[\Mmat]$.

\begin{proposition}[Expected exposure operator]\label{prop:mean-exposure}
\begin{equation}\label{eq:mean-Q}
  \bar{\Qmat}\equiv\mathbb E[\Qmat]=\Jmat^T\,\bar\Mmat\,\Jmat,
  \qquad
  \mathbb E[P_\mathrm{abs}(\xx)]=\xx^H\bar{\Qmat}\,\xx.
\end{equation}
\end{proposition}

\begin{proof}
$\Qmat=\Jmat^T\Mmat\Jmat$ is linear in $\Mmat$ for fixed $\Jmat$, and
$\Jmat$ is deterministic (combinatorial).  Linearity of expectation gives
$\bar\Qmat=\Jmat^T\bar\Mmat\Jmat$.  Then
$\mathbb E[\xx^H\Qmat\xx]=\xx^H\bar\Qmat\xx$ for deterministic $\xx$.
\end{proof}

The expected-value ECBF problem is
\begin{equation}\label{eq:mean-ECBF}
  \max_\xx|\hh^T\xx|^2\;\text{s.t.}\;\xx^H\bar\Qmat\xx\le P_\mathrm{abs}^{\max},\;\|\xx\|^2\le P,
\end{equation}
and inherits the closed-form solution of \Cref{thm:x-path} with
$\Mmat\mapsto\bar\Mmat$.  No additional machinery is needed.

\begin{remark}[Robust ECBF]\label{rem:robust}
When one wishes to control $\Pr[P_\mathrm{abs}>P_\mathrm{abs}^{\max}]\le\epsilon$ rather than
the mean, the quadratic form $\xx^H\Qmat\xx$ is random.  For a Gaussian
channel $\psin\sim\mathcal{CN}$ with path-covariance $\mathbf C$, the
moment-generating function $\mathbb E[e^{t\xx^H\Qmat\xx}]$ is the
Gaussian quadratic MGF $\det(\Id-2t\,\mathbf C^{1/2}\xx\xx^H\mathbf C^{1/2})^{-1}$,
yielding Chernoff bounds of the form
$\Pr[P_\mathrm{abs}>P_\mathrm{abs}^{\max}]\le\exp\!\bigl(-t P_\mathrm{abs}^{\max}+\Psi(t,\xx)\bigr)$.
Inner minimisation over $t$ and outer maximisation in $\xx$ define a
robust ECBF; the path-space factorisation reduces the matrix inverse to
the same $N\times N$ kernel $(\nu\Id_N+\lambda\bar\Mmat\Pmat)$ as in
\Cref{thm:x-path}, pre-shaped by $\mathbf C$.
\end{remark}

\begin{remark}[Ergodic trace]\label{rem:ergodic-tr}
Combining \Cref{prop:mean-exposure} with \Cref{prop:iso},
\[
  \mathbb E_{\xx,\mathrm{paths}}[P_\mathrm{abs}]
  =\frac{P}{M}\operatorname{tr}(\bar\Qmat)=\frac{P}{M}\operatorname{tr}(\bar\Mmat\,\Pmat),
\]
i.e.\ under joint isotropy of precoder and path ensemble, the absorbed
power is a single scalar: the trace of the mean path-Gram restricted to
the element-block subspace, per transmit degree of freedom.
\end{remark}

% ============================================================
\section{Low-rank body approximation}\label{sec:lowrank}
% ============================================================

In practice the body behaves as a low-rank scatterer with respect to a
fixed path set: only a handful of absorption modes carry appreciable
weight.  This is already implicit in the eigendecomposition of $\Qmat$
(\cref{sec:exposure-operator} of the monograph) but becomes an
exploitable computational structure in path space.

\begin{definition}[Rank-$k$ path-space truncation]\label{def:trunc}
Let $\Qmat=\sum_{k=1}^M\mu_k\vv_k\vv_k^H$ be the eigendecomposition of the
element-space exposure operator, with $\mu_1\ge\cdots\ge\mu_M\ge 0$ and
$\vv_k\in\Complex^M$ orthonormal.  The \emph{rank-$r$ approximation} is
\begin{equation}\label{eq:Q-trunc}
  \Qmat^{(r)}\;=\;\sum_{k=1}^{r}\mu_k\,\vv_k\vv_k^H,\qquad 1\le r\le M.
\end{equation}
\end{definition}

\begin{proposition}[Approximation error]\label{prop:trunc-err}
For every precoder with $\|\xx\|^2\le P$,
\begin{equation}\label{eq:trunc-err}
  \bigl|\xx^H(\Qmat-\Qmat^{(r)})\xx\bigr|\;\le\;\mu_{r+1}\,P,
\end{equation}
with equality when $\xx$ is aligned to $\vv_{r+1}$.  Consequently the
exposure constraint $\xx^H\Qmat^{(r)}\xx\le P_\mathrm{abs}^{\max}-\mu_{r+1}P$
is a sufficient condition for $\xx^H\Qmat\xx\le P_\mathrm{abs}^{\max}$.
\end{proposition}

\begin{proof}
$(\Qmat-\Qmat^{(r)})=\sum_{k>r}\mu_k\vv_k\vv_k^H$ is Hermitian PSD with
operator norm $\mu_{r+1}$.  Hence
$\xx^H(\Qmat-\Qmat^{(r)})\xx\le\mu_{r+1}\|\xx\|^2\le\mu_{r+1}P$.
\end{proof}

\begin{theorem}[Path-space low-rank equivalence]\label{thm:lowrank-path}
The top-$r$ eigenvalues of $\Qmat$ coincide with the top-$r$ nonzero
eigenvalues of $\Mmat\Pmat$.  Moreover, if $\Mmat\Pmat\,\uu_k=\mu_k\uu_k$
with $\mu_k>0$, then $\vv_k\propto\Jmat^T\uu_k$.  Therefore
\begin{equation}\label{eq:Q-trunc-path}
  \Qmat^{(r)}\;=\;\sum_{k=1}^{r}\mu_k\,
  \frac{\Jmat^T\uu_k\uu_k^H\Jmat}{\uu_k^H\Pmat\uu_k},
\end{equation}
and a rank-$r$ exposure-aware precoder can be computed using $r$
path-space eigenpairs at cost $O(r\cdot\|\Mmat\|_0)$ per Lanczos
iteration, where $\|\Mmat\|_0$ is the number of non-negligible entries of
$\Mmat$.
\end{theorem}

\begin{proof}
By \Cref{cor:lam-path} the nonzero spectra match.  Normalisation: from
$\Qmat\,\Jmat^T\uu_k=\Jmat^T\Mmat\Jmat\,\Jmat^T\uu_k=\Jmat^T(\Mmat\Pmat)\uu_k
=\mu_k\Jmat^T\uu_k$ (since $\uu_k$ is an eigenvector of $\Mmat\Pmat$),
$\Jmat^T\uu_k$ is an (unnormalised) eigenvector of $\Qmat$ at eigenvalue
$\mu_k$.  Its squared element-space norm is
$\|\Jmat^T\uu_k\|^2=\uu_k^H\Jmat\Jmat^T\uu_k=\uu_k^H\Pmat\uu_k$, giving
\eqref{eq:Q-trunc-path}.  The Lanczos cost is standard.
\end{proof}

\begin{remark}[Dosimetric SVD]\label{rem:svd}
\Cref{thm:lowrank-path} realises the ``dosimetric SVD'' idea: the
eigenpairs $\{(\mu_k,\uu_k)\}$ are body modes in path space.  They replace
the angular directivity basis $D(\khat)$ of Part~I with a body-mode basis
adapted to the particular path set.  Truncation at $r$ controls the
worst-case absorbed-power error by $\mu_{r+1}P$ via \Cref{prop:trunc-err},
independent of $\|\xx\|$.
\end{remark}

% ============================================================
\section{ECBF as a generalised eigenvalue problem}\label{sec:gep}
% ============================================================

The closed-form ECBF solution of \Cref{thm:x-path} follows from KKT with
two multipliers.  A complementary viewpoint drops the power constraint
and asks for the precoder that maximises signal per unit absorbed power;
this is a single generalised eigenvalue problem (GEP).

\begin{proposition}[Signal-to-exposure GEP]\label{prop:GEP}
Assume $\Qmat\succ0$ (generic full-rank body).  Then
\begin{equation}\label{eq:GEP}
  \max_{\xx\ne 0}\;\frac{|\hh^T\xx|^2}{\xx^H\Qmat\xx}
  \;=\;\hh^T\Qmat^{-1}\hh^*,
\end{equation}
attained at $\xx^\star\propto\Qmat^{-1}\hh^*$.  In path space,
\begin{equation}\label{eq:GEP-path}
  \max_{\xx\ne 0}\;\frac{|\hh^T\xx|^2}{\xx^H\Qmat\xx}
  \;=\;\bvec^H\,\Jmat(\Jmat^T\Mmat\Jmat)^{-1}\Jmat^T\,\bvec.
\end{equation}
\end{proposition}

\begin{proof}
The objective is the Rayleigh quotient of the pencil $(\hh^*\hh^T,\Qmat)$
with $\hh^*\hh^T$ rank-one Hermitian PSD.  Generalised eigenproblem:
$\hh^*\hh^T\xx=\mu\Qmat\xx$.  Left-multiplying by $\hh^T\Qmat^{-1}$ yields
$\hh^T\Qmat^{-1}\hh^*\cdot\hh^T\xx=\mu\,\hh^T\xx$, so for any eigenvector
with $\hh^T\xx\ne 0$, $\mu=\hh^T\Qmat^{-1}\hh^*$.  Substituting $\xx^\star
=\Qmat^{-1}\hh^*$ into the Rayleigh quotient verifies the maximum and its
value.  The path-space form follows from $\hh^*=\Jmat^T\bvec$,
$\hh^T=\bvec^H\Jmat$, and $\Qmat=\Jmat^T\Mmat\Jmat$.
\end{proof}

\begin{remark}[ECBF in the GEP picture]\label{rem:gep-ecbf}
The ECBF precoder of \Cref{thm:x-path} is the dual of this GEP under the
added power constraint.  As $\lambda\to 0$ (exposure constraint slack),
$\xx^\star\to\hh^*$ (MRT, maximises signal ignoring body).  As
$\nu\to 0$ (power constraint slack), $\xx^\star\to\Qmat^{-1}\hh^*$ (GEP
optimum, maximises signal per unit absorbed power).  The Pareto frontier
between these limits traces the ECBF family indexed by $\lambda/\nu$.
\end{remark}

\begin{remark}[Geometric reading in path space]\label{rem:oblique}
Let $\mathbf y\equiv\Jmat\xx\in\mathcal E$ be the path-space lift of the
precoder.  The GEP \eqref{eq:GEP} rewrites as
\begin{equation}\label{eq:GEP-y}
  \max_{\mathbf y\in\mathcal E\setminus\{0\}}\;
  \frac{|\bvec^H\mathbf y|^2}{\mathbf y^H\Mmat\mathbf y},
\end{equation}
a constrained Rayleigh quotient.  On unconstrained $\Complex^N$ the
maximiser would be $\mathbf y=\Mmat^{-1}\bvec$ with value
$\bvec^H\Mmat^{-1}\bvec$; constraining $\mathbf y\in\mathcal E$ replaces
it by $\mathbf y^\star=\Pi_\mathcal E^{\Mmat}\Mmat^{-1}\bvec
=\Jmat(\Jmat^T\Mmat\Jmat)^{-1}\Jmat^T\bvec$, the $\Mmat$-orthogonal
projection of the unconstrained optimum onto the element-tied subspace.
Here $\Pi_\mathcal E^{\Mmat}\equiv\Jmat(\Jmat^T\Mmat\Jmat)^{-1}\Jmat^T\Mmat$
is the $\Mmat$-orthogonal projector onto $\mathcal E$; it is idempotent
and self-adjoint with respect to $\langle u,v\rangle_{\Mmat}=u^H\Mmat v$.
The GEP value in \eqref{eq:GEP-path} is then the squared
$\Mmat$-inner-product length of $\mathbf y^\star$.  The price of the
element-tied constraint is the loss
$\bvec^H\Mmat^{-1}\bvec-\bvec^H\Jmat\Qmat^{-1}\Jmat^T\bvec\ge 0$: the
signal-per-absorbed-power gap between a hypothetical
path-by-path precoder and the physical per-element precoder.
\end{remark}

% ============================================================
\section{Reciprocity and the body as a virtual array}\label{sec:reciprocity}
% ============================================================

A short observation that illuminates $\Qmat$ physically.  Consider the
adjoint situation: a dipole source placed at body point $\rr\in\Sigmabody$
with polarisation $\hat{\mathbf u}\in\Complex^3$ radiates into free
space; the monochromatic field at BS element $j$ has some complex
amplitude $\mathcal G_{j,\hat{\mathbf u}}(\rr)$.  Electromagnetic
reciprocity (specifically Lorentz reciprocity for line-of-sight Fresnel
transmission) identifies $\mathcal G_{j,\hat{\mathbf u}}(\rr)$ with the
$(j,\hat{\mathbf u})$ component of
$\tilde{\mathbf g}_j(\rr)=(\Gtil(\rr)\,\mathbf e_j)$ in the exposure
branch, up to tissue-absorption conventions that do not affect
proportionality.

Under this identification, the exposure operator $\Qmat$ is the transmit
Gram of the ``virtual body array'' whose aperture is the integral over
$\Sigmabody$:
\begin{equation}\label{eq:reciprocity}
  \Qmat\;=\;\int_\Sigmabody\sum_{\hat{\mathbf u}\in\{\hat{\mathbf u}_1,\hat{\mathbf u}_2,\hat{\mathbf u}_3\}}
  \mathcal G_{\cdot,\hat{\mathbf u}}(\rr)\,\mathcal G_{\cdot,\hat{\mathbf u}}(\rr)^H\,\diff A,
\end{equation}
with $\{\hat{\mathbf u}_k\}$ any orthonormal basis of $\Complex^3$.  The
top eigenvector $\vv_1$ of $\Qmat$ is then the MRT precoder for the best
coupling to \emph{any} body-distributed virtual source, and
$\lambda_{\max}(\Qmat)$ is the squared maximum coupling gain.  The
signal branch $\mathbf h$ is the scalar version of the same construction
for a single UE-located source.  Signal and exposure are then two
reciprocal MRT problems sharing the path geometry $(\Jmat,\Mmat)$ and
differing only in whether the virtual source is a single point
($\rr_\mathrm{UE}$) or an integrated aperture ($\Sigmabody$).

This is the physical content of the Rayleigh quotient for $\rho$: it
measures how well the single-point MRT precoder aligns with the
distributed-aperture MRT precoder.

% ============================================================
\section{Where this formalism points}\label{sec:outlook}
% ============================================================

With the formalism in place, several directions are now concrete
problems rather than open metaphors.

\paragraph{Fourier-on-$\Sigmabody$ calculus for canonical phantoms.}
\Cref{rem:fourier} shows that $\Mmat$ is a surface Fourier transform of a
path-pair kernel evaluated at wave-vector differences
$\mathbf q_{nn'}=k_0(\khat_{n'}-\khat_n)$.  Ellipsoids, capsules, and
cylinders have analytically known surface Fourier spectra; coupling these
with per-path Fresnel factors gives closed-form $\Mmat$ for canonical
phantoms, the Part~III counterpart of the Part~I absorption directivity
$D(\khat)$.  The near-field version (\Cref{rem:confocal}) needs the same
calculation with confocal coordinates in place of plane waves.

\paragraph{Measurement-driven $\Mmat$.}  The path-exposure Gram
$\Mmat\in\Complex^{N\times N}$ has $N^2$ entries; many may be dominated
by a low-rank or sparse structure (\Cref{thm:lowrank-path}).  A sparse
sensing problem: probe the body with $r\ll N$ calibrated single-path
excitations (measurable via per-element pattern calibrations) and recover
$\Mmat$ by matrix completion, using the eigendecomposition of $\Qmat$ as
an anchor.  The estimated $\Mmat$ then feeds an online ECBF.

\paragraph{Time-dynamic exposure.}  When the body moves, $\Mmat$ evolves
continuously.  First-order perturbation theory on the PSD operator
$\Mmat\Pmat$ predicts how the dominant modes $(\mu_k,\uu_k)$ drift with
pose; tracking these locally-linear updates is cheaper than rebuilding
$\Qmat$ per frame.  \Cref{rem:ergodic-tr} gives the corresponding
ergodic statistic.

\paragraph{Higher-order corrections.}  Approximations~1 and~2 of
Part~III enter $\Mmat$ only through per-pair entries; curvature
(\cref{sec:curvature} of the monograph) and diffraction
(\cref{sec:diffraction}) corrections act as perturbations
$\Mmat\to\Mmat+\Delta\Mmat$ with $\|\Delta\Mmat\|$ bounded by the Part~III
error budget.  The GEP and push-through remain valid under such
perturbations by spectral stability (Weyl, Davis--Kahan), turning the
Part~III error bounds into quantitative robustness statements for ECBF.

\paragraph{Coupled multi-band MU-MIMO.}  Combining
\Cref{prop:omega-pt,prop:mean-exposure} gives a broadband, stochastic,
multi-user exposure-constrained design in which each frequency bin
contributes an independent path-space problem, coupled only through
integrated dual multipliers.  Water-filling-like allocation with
exposure ``costs'' $\mu_k(\omega)$ is the natural next step.

\paragraph{Summary.}  The entire Part~III pipeline, together with its
broadband, near-field, stochastic, low-rank, and multi-user extensions,
collapses onto one chain of matrix factorisations,
\[
  \xx\;\xrightarrow{\;\Jmat\;}\;\Complex^N
  \;\xrightarrow{\;\Phimat(\rr)\;}\;\Complex^N
  \;\xrightarrow{\;\Psimat,\alphavec^T,\Psitil(\rr)\;}\;
  \text{field, signal, exposure.}
\]
The path-exposure Gram
$\Mmat=\int_\Sigmabody\Phimat^H\Psitil^H\Psitil\Phimat\,\diff A$ is the
single $N\times N$ object encoding body response; the dispatch $\Jmat$
is the bridge between the $M$-dimensional control space and the
$N$-dimensional path space; the element-tied subspace
$\mathcal E=\operatorname{range}(\Jmat)$ is the geometric locus of
reachable absorption modes.  The monograph's central quantities are
pre-diagonalised by this chain:
\[
  \Qmat=\Jmat^T\Mmat\Jmat,\qquad
  \lambda_{\max}(\Qmat)=\lambda_{\max}(\Mmat\Pmat),\qquad
  \xx^\star\propto\Jmat^T(\nu\Id_N+\lambda\Mmat\Pmat)^{-1}\bvec.
\]
Each of the six generalisations of this note
(\Cref{sec:trace,sec:broadband,sec:nearfield,sec:stochastic,sec:lowrank,sec:gep})
enters as a small perturbation or structural variation of $\Mmat$; none
requires new matrix machinery.

\end{document}
